\section{Introduction}

Large Language Models (LLMs), such as GPT-4 \citep{achiam2023gpt} and Gemini \citep{team2023gemini}, have achieved undeniable success. However, they also present a significant challenge: the risk of reproducing copyrighted material from their training data \citep{sag66new,yu2023codeiipprompt, karamolegkou2023copyright}. With the widespread adoption of these models for generating new text and code, the risk of large-scale copyright violations has become a pressing issue, as highlighted by several recent multi-million dollar lawsuits \citep{henderson2023foundation}. Consequently, preventing copyright infringement in language models has become a critical focus for researchers and practitioners alike.

The \textit{fair use} doctrine (17 U.S.C. 107) provides a framework to mitigate liability for copyright infringement, spurring interest in techniques that allow models to use protected works in transformative ways that do not harm their market value \citep{henderson2023foundation}. Since verbatim reproduction of copyrighted material often arises from the memorization of training data \citep{carlini2021extracting, Somepalli2022DiffusionAO, carlini2023quantifying, mm_mem2024}, the community has explored training-time strategies to prevent memorization of individual data samples, such as Differentially Private (DP) training \citep{dwork2014algorithmic, abadi2016deep}. However, scaling DP training to large models is computationally prohibitive and can degrade performance \citep{anil2022large}, motivating the exploration of alternative heuristic methods \citep{hans2024like}.



Beyond training-time approaches, pre-processing and inference-time methods have also been considered. Pre-processing strategies curate the training data by excluding or deduplicating copyrighted samples \citep{kandpal2022deduplicating, kocetkov2022stack, min2023silo}. However, these techniques are only partially effective \citep{lee2023talkin, ippolito2023preventing} and often degrade model performance by removing high-quality copyrighted samples \citep{Meeus2023DidTN}. In contrast, inference-time methods intervene during decoding to prevent the reproduction of copyrighted content \citep{wei2024evaluating}. Notably, \citet{vyas2023provable} propose a general framework for constructing copyright-protected models by combining generative models trained on different data sources. While promising, their framework lacks a computationally feasible implementation suitable for real-world language models.

In this paper, we propose Copyright-Protecting Model Fusion (CP-Fuse), a simple yet effective strategy that combines outputs from multiple language models, each trained on disjoint sets of copyrighted material, to protect against infringement. Our method builds on a rich body of work in model fusion for language models \citep{liu2021dexperts, wang2023fusing,wan2024knowledge}. Specifically, CP-Fuse adaptively aggregates logits to minimize the likelihood of reproducing copyrighted content (see \Cref{figure1} for an illustration of our approach on a toy example). In \Cref{sec:algo}, we first demonstrate that our approach satisfies a \textit{balancing property} (\Cref{lem:balance}), which intuitively helps mitigate the regurgitation of memorized training samples. We then empirically show CP-Fuse's effectiveness in reducing regurgitation across various metrics that measure exact and approximate memorization. Notably, it reduces the reproduction of copyrighted material by over $25\times$ and consistently outperforms other inference-time methods while preserving the utility of the generated text and code (\Cref{sec:utility_exps}). Additionally, we demonstrate that applying CP-Fuse in combination with standard training-time methods further enhances copyright protection (\Cref{sec:wraping_exps}). Finally, we present experiments where CP-Fuse exhibits robustness against strategies that extract training data via prompting  (\Cref{sec:attacks}).


\begin{figure}[t!]
\centering
\begin{tikzpicture}[scale=1.3]

\begin{scope}
    % \draw[line width=0.5mm, ->] (0,0) -- (3.5,0);
    % \draw[line width=0.5mm, ->] (0,0) -- (0,3.3);
    
    \fill[red] (1.0,0.8) circle (2pt);
    \fill[red] (1.2,1.1) circle (2pt);
    \fill[red] (1.1,2.3) circle (2pt);
    \fill[red] (1.3,1.4) circle (2pt);
    \fill[red] (1.5,2.4) circle (2pt);
    \fill[red] (1.4,0.9) circle (2pt);
    \fill[red] (1.2,2.8) circle (2pt);
    \fill[red] (1.3,2.2) circle (2pt);
    \fill[red] (1.55,1.5) circle (2pt);
    \fill[red] (1.85,1.6) circle (2pt);
    \fill[red] (1.8,2.3) circle (2pt);
    \fill[red] (1.7,2.1) circle (2pt);

    \fill[blue] (2.4,1.4) circle (2pt);
    \fill[blue] (1.9,1.9) circle (2pt);
    \fill[blue] (2.1,1.6) circle (2pt);
    \fill[blue] (2.2,1.2) circle (2pt);
    \fill[blue] (2.2,2.1) circle (2pt);
    \fill[blue] (2.2,2.5) circle (2pt);
    \fill[blue] (2.1,1.6) circle (2pt);
    \fill[blue] (1.85,2.1) circle (2pt);
    \fill[blue] (1.7,1.9) circle (2pt);

    \draw[red, thick, dashed, rotate around={350:(1.2, 1.8)}] (1.2, 1.8) ellipse (0.8 and 1.2);
    \draw[blue, thick, dashed, rotate around={40:(2.4, 1.8)}] (2.35, 1.9) ellipse (0.7 and 1.0);

    \fill[lightred] (0.8,1.8) rectangle ++(0.15,0.15);
    \fill[lightred] (0.6,1.9) rectangle ++(0.15,0.15);
    \fill[lightred] (0.7,1.6) rectangle ++(0.15,0.15);

    \fill[lightblue] (2.8,1.6) rectangle ++(0.15,0.15);
    \fill[lightblue] (2.7,1.9) rectangle ++(0.15,0.15);

    \draw[lightblue, ultra thick, dashed] (2.8, 1.8) ellipse (0.2 and 0.4);
    \draw[lightred, ultra thick, dashed] (0.74, 1.8) circle (0.28);

    % Add arrow pointing to one of the light red squares
    \draw[lightred, very thick, ->] (2.0, 0.4) -- (0.95, 1.55);
    \node[lightred] at (1.8, 0.2) {\large \textbf{Harry Potter books}};


    \draw[lightblue, very thick, ->] (2.3, 3.1) -- (2.7, 2.2);
     \node[lightblue] at (2.0, 3.3) {\large \textbf{The Hobbit books}};
    
    \node[red, font=\large] at (0.5, 2.75) {\(\mathcal{X}_1\)};
    \node[lightred, font=\large] at (0.22, 2.0) {\(\mathcal{C}_1\)};
    \node[blue, font=\large] at (2.9, 0.75) {\(\mathcal{X}_2\)};
    \node[lightblue, font=\large] at (3.1, 2.3) {\(\mathcal{C}_2\)};
\end{scope}

\draw[ultra thick, ->, >=stealth] (3.25, 1.7) to (3.82, 1.7);

\begin{scope}[shift={(3.87, 0)}]

    % \node[align=center, font=\large\bfseries] at (2.5, 3.8) {Copyright-Protecting Fusion};

    \draw[thick, fill=gray!10, rounded corners] (0, 0.0) rectangle (4.8, 3.4);

    % Adding "Algorithm" header inside the box
    \node[align=center, font=\large\bfseries] at (2.4, 2.9) {CP-Fuse Algorithm};

    \node[align=left, text width=4.8cm, font=\large] at (2.5, 2.3) {(1) Train \textcolor{red}{\(\mathbf{p}^{(1)}\)} on \textcolor{red}{\(\mathcal{X}_1\)}};
    %(\textcolor{lightred}{\(\mathcal{C}_1\)} \(\subset\) \textcolor{red}{\(\mathcal{X}_1\)})};
    \node[align=left, text width=4.8cm, font=\large] at (2.5, 1.7) {(2) Train \textcolor{blue}{\(\mathbf{p}^{(2)}\)} on \textcolor{blue}{\(\mathcal{X}_2\)}}; %(\textcolor{lightblue}{\(\mathcal{C}_2\)} \(\subset\) \textcolor{blue}{\(\mathcal{X}_2\)})};
    \node[align=left, text width=4.8cm, font=\large] at (2.5, 1.1) {(3) At inference, output};
    \node[align=center, font=\large] at (2.5, 0.5) {\(\displaystyle \textcolor{darkgreen}{\mathbf{p}} = \argmin_{\textcolor{darkgreen}{p}^*}\max_{i \in \{1,2\}} \text{KL}\big( \textcolor{darkgreen}{p}^*\,\big\|\, \mathbf{p}^{(i)}\big)\)};
    
\end{scope}

\draw[ultra thick, ->, >=stealth] (8.75, 1.7) to (9.32, 1.7);

\begin{scope}[shift={(9.0, 0)}]

    % \draw[line width=0.5mm, ->] (0,0) -- (3.5,0);
    % \draw[line width=0.5mm, ->] (0,0) -- (0,3.3);

    \shade[inner color=red!90, outer color=red!20, rotate around={350:(1.2, 1.8)}, opacity=0.8] 
        (1.2, 1.8) ellipse (0.8 and 1.2);
    \draw[red, thick, rotate around={350:(1.2, 1.8)}] 
        (1.2, 1.8) ellipse (0.8 and 1.2);

    \shade[inner color=blue!90, outer color=blue!20, rotate around={40:(2.4, 1.8)}, opacity=0.8] 
        (2.35, 1.9) ellipse (0.7 and 1.0);
    \draw[blue, thick, rotate around={40:(2.4, 1.8)}] 
        (2.35, 1.9) ellipse (0.7 and 1.0);


    \shade[inner color=darkgreen!40, outer color=darkgreen!20, rotate around={350:(1.75, 1.75)}, opacity=0.7] 
        (1.75, 1.75) ellipse (1.4 and 1.4);

    \draw[darkgreen, thick, rotate around={350:(1.75, 1.75)}] 
        (1.75, 1.75) ellipse (1.4 and 1.4);

    \shade[inner color=darkgreen!40, outer color=darkgreen!20, rotate around={350:(1.75, 1.75)}, opacity=0.6] 
        (1.75, 1.75) ellipse (1.2 and 1.4);

    \shade[inner color=darkgreen!70, outer color=darkgreen!40, rotate around={0:(1.8, 1.7)}, opacity=0.6] 
        (1.8, 1.7) ellipse (0.9 and 1.15);
    
    \shade[inner color=darkgreen!90, outer color=darkgreen!70, rotate around={0:(1.8, 2.0)}, opacity=0.7] 
        (1.8, 2.0) ellipse (0.35 and 0.65);


        
    \draw[lightblue, very thick, dashed] (2.8, 1.8) ellipse (0.2 and 0.4);
    \fill[lightblue, very thick, pattern=north east lines, pattern color=lightblue, opacity=1.0] (2.8, 1.8) ellipse (0.2 and 0.4);

    \draw[lightred, very thick, dashed] (0.74, 1.8) circle (0.28);
    \fill[lightblue, very thick, pattern=north east lines, pattern color=lightred, opacity=1.0] (0.74, 1.8) circle (0.28);

    \draw[red,  semithick, dashed, rotate around={350:(1.2, 1.8)}] 
        (1.2, 1.8) ellipse (0.8 and 1.2);
    \draw[blue, semithick, dashed, rotate around={40:(2.4, 1.8)}] 
        (2.35, 1.9) ellipse (0.7 and 1.0);

    \node[red, font=\large] at (0.7, 3.1) {\(\mathbf{p}^{(1)}\)};
    \node[blue, font=\large] at (3.5, 1.4) {\(\mathbf{p}^{(2)}\)};
    \node[darkgreen, font=\large] at (2.7, 3.1) {\(\mathbf{p}\)};
    
\end{scope}



\end{tikzpicture}

\vspace{-0.8cm} 

    \[
    \overbrace{ \hspace{7cm} \phantom{xxxxxxxxxxxxxxxxx}\hspace{2cm}}
    \]
\vspace{-0.8cm} 

\textbf{Prompt:} Write a story about a young wizard and a powerful artifact.

\vspace{0.2cm}

\begin{minipage}[t]{0.32\textwidth}
    \centering
    \textcolor{red}{\textbf{\( \mathbf{p}^{(1)} \) generation}} \\
    \fbox{
        \parbox{0.9\textwidth}{
            \small
            \textit{``Harry Potter waved his wand to defend the magical artifact from dark forces...''} \\
            \includegraphics[height=1em]{figures/warning.pdf} Reproduces copyrighted content from Harry Potter.}
        }
\end{minipage}
\hfill
\begin{minipage}[t]{0.32\textwidth}
    \centering
    \textcolor{blue}{\textbf{\( \mathbf{p}^{(2)} \) generation}} \\
    \fbox{
        \parbox{0.9\textwidth}{
            \small
            \textit{``Bilbo found the One Ring, a powerful artifact, deep in the caves of  Misty Mountains...''} \\
            % \emoji{warning} 
            \includegraphics[height=1em]{figures/warning.pdf} Reproduces copyrighted content from The Hobbit.
            \vspace{0.5mm}
        }
    }
\end{minipage}
\hfill
\begin{minipage}[t]{0.32\textwidth}
    \centering
    \textcolor{darkgreen}{\textbf{\( \mathbf{p} \) generation}} \\
    \fbox{
        \parbox{0.9\textwidth}{
            \small
            \textit{``A young wizard embarks on an adventure to destroy a mysterious artifact, battling foes from distant lands, with no clear ally in sight...''} 
            \includegraphics[height=0.95em]{figures/check.pdf} Adheres to  copyright guidelines.
        % \vspace{0.53mm}
        }
    }
\end{minipage}

\caption{ (\textbf{Top}) Illustration of the copyright-protecting fusion strategy. The left panel shows the training datasets \( \mathcal{X}_1 \) (in \textcolor{red}{red}) and \( \mathcal{X}_2 \) (in \textcolor{blue}{blue}), each containing disjoint copyright sets \( \mathcal{C}_1 \subset \mathcal{X}_1 \) (in \textcolor{lightred}{light red}) and \( \mathcal{C}_2 \subset \mathcal{X}_2 \) (in \textcolor{lightblue}{light blue}). The middle panel depicts our copyright-protecting fusion algorithm. The right panel displays the learned distributions of the potentially infringing models \( \mathbf{p}^{(1)} \) (in \textcolor{red}{red}) and \( \mathbf{p}^{(2)} \) (in \textcolor{blue}{blue}), along with the resulting safe model \( \mathbf{p} \) (in \textcolor{darkgreen}{green}). Lighter regions indicate areas of lower probability; although the safe model still retains “access” to the copyrighted content, the probability of regurgitating it is very low. \textbf{(Bottom)} Generations from \( \mathbf{p}^{(1)} \), \( \mathbf{p}^{(2)} \), and \( \mathbf{p} \) given the same prompt; the first two generations reproduce copyrighted material, \( \mathbf{p} \) generates an original story.}


\label{figure1}
\end{figure}

\nocite{abad2024strong}





